\section{Related Work}
\label{sec:related_work}

\subsection{LLM Serving Systems and Memory Management}
\label{subsec:rw-serving}

Orca~\cite{yu2022orca} introduced iteration-level scheduling for
transformer-based generative models, allowing individual requests to enter
and leave a batch between decoding steps rather than waiting for the entire
batch to complete, which is the continuous-batching mechanism this work
depends on. vLLM~\cite{kwon2023vllm} builds on this idea and additionally
introduces PagedAttention, which manages the key-value cache in fixed-size
blocks to reduce memory fragmentation and increase the number of sequences
that can be served concurrently on a fixed amount of GPU memory. Both systems
focus on maximizing GPU utilization and aggregate throughput through
memory-efficient batching; neither prescribes a specific admission-ordering
policy for the requests that populate each batch, which is the aspect of the
serving stack this work targets.

\subsection{Output-Length-Aware and Learning-to-Rank Scheduling}
\label{subsec:rw-scheduling}

A separate line of work argues that batching efficiency alone is insufficient
if the admission order itself causes short requests to wait behind long ones,
and proposes estimating output length ahead of time to inform scheduling.
Response Length Perception and Sequence Scheduling~\cite{zheng2023responselength}
frames this as a classification problem, training a model to bucket prompts
into coarse output-length classes and prioritizing shorter buckets; the
classification-based scheduler evaluated in this paper follows that design.
Efficient LLM Scheduling by Learning to Rank~\cite{fu2024efficient}, the work
this project directly extends, instead trains a compact OPT-125M predictor
with a listwise ranking loss so that its output score orders requests
consistently with their true relative output lengths, without requiring the
predictor to commit to a specific length bucket or exact token count. The
underlying ranking objective follows the learning-to-rank literature more
broadly, including pairwise formulations such as RankNet~\cite{burges2005ranknet}
and listwise formulations such as ListMLE~\cite{xia2008listmle}, the loss used
to train the original predictor in~\cite{fu2024efficient}. This project keeps
that learned predictor intact and asks whether combining its score with an
independent, non-learned signal -- prompt length -- improves robustness when
the predictor's score alone is imperfect, which to our knowledge has not been
examined in prior scheduling work.
